\section{Algoritmi di ricerca}

In questa sezione viene descritta la ricerca di base adottata per il
problema del giro del corriere. La procedura utilizza una visita in
profondità (DFS) per determinare l'ordine di esplorazione degli stati,
integrata con un meccanismo di branch-and-bound per scartare i rami
che non possono migliorare la migliore soluzione completa già trovata.

\subsection{Rappresentazione dello stato e spazio di ricerca}

Prima di descrivere l'algoritmo, si definiscono i concetti
di stato e spazio di ricerca.

Uno \emph{stato} rappresenta una configurazione del problema
in un dato momento della ricerca. Nel problema del giro del
corriere, ciascuno stato contiene il percorso parziale
costruito fino a quel momento, l'insieme delle consegne
ancora da effettuare e la distanza accumulata.
La posizione attuale del corriere corrisponde all'ultimo
punto del percorso.

Lo stato iniziale contiene il solo deposito nel percorso,
tutte le consegne ancora da effettuare e un costo pari a zero.
Quando tutte le consegne sono state visitate, il giro viene
completato tornando al deposito. La distanza di ritorno
viene inclusa nel costo finale.

Lo \emph{spazio di ricerca} è l'insieme degli stati
raggiungibili dallo stato iniziale attraverso le azioni
consentite. In questo caso, comprende i percorsi ammissibili che
partono dal deposito e visitano un sottoinsieme delle consegne,
senza ripetizioni. Gli stati che hanno visitato tutte le consegne
costituiscono le foglie dell'albero di ricerca; a ciascuna foglia
si associa un giro completo aggiungendo il ritorno al deposito.


\subsection{Espansione e generazione dei successori}

L'espansione di uno stato consiste nella generazione dei possibili
proseguimenti del percorso corrente.

Sia $V = \{d, c_1, \ldots, c_n\}$ l'insieme delle località,
dove $d$ rappresenta il deposito e $c_1, \ldots, c_n$
i punti di consegna. Per ciascuno stato si distinguono
l'insieme $C \subseteq V$ delle località già visitate,
comprendente il deposito, e l'insieme
$N = V \setminus C$ delle consegne ancora da effettuare.
Si ha quindi:
\[
    C \cap N = \emptyset,
    \qquad
    C \cup N = V.
\]

Lo stato viene indicato con $s=(p,N,g)$, dove $p$ conserva
l'ordine delle località visitate e $g$ indica il costo accumulato.
L'insieme $C$ si ricava dal percorso $p$ e non deve quindi essere
memorizzato separatamente.

Per ogni consegna $v \in N$ viene generato uno stato
successore. Il nuovo percorso si ottiene aggiungendo $v$
alla fine di $p$, mentre gli insiemi vengono aggiornati come:
\[
    C' = C \cup \{v\},
    \qquad
    N' = N \setminus \{v\}.
\]
Indicando con $u$ l'ultima località del percorso corrente,
il costo del successore è:
\[
    g' = g + \operatorname{dist}(u,v).
\]
Ogni successore rappresenta un'alternativa distinta: gli aggiornamenti
non modificano lo stato padre né gli altri successori.

\subsection{Ordine di visita tramite DFS}

Per determinare l'ordine di esplorazione degli stati, l'algoritmo utilizza
una \textbf{visita in profondità} (DFS, \emph{Depth-First Search}),
implementata mediante una pila con politica LIFO: l'ultimo stato inserito
è il primo a essere estratto. La DFS fornisce una procedura di base
semplice da descrivere e da integrare con la potatura.
L'Algoritmo~\ref{alg:dfs} presenta esclusivamente la strategia di visita;
il criterio di branch-and-bound viene descritto nella sottosezione
successiva.

\begin{algorithm}[htbp]
\caption{Depth-First Search}
\label{alg:dfs}
\begin{algorithmic}[1]

\Require Stato iniziale $s_0$
\Ensure Visita di tutti gli stati dell'albero in profondità

\State $S \gets$ pila vuota \Comment{Conserva gli stati in attesa}
\State $\operatorname{push}(S,s_0)$

\While{$S$ non è vuota}
    \State $s \gets \operatorname{pop}(S)$ \Comment{Estrae l'ultimo inserito}
    \State \textbf{visita} $s$
    \State $L \gets \operatorname{Successori}(s)$ \Comment{Genera le estensioni ammissibili}
    \State Ordina $L$ per indice crescente della consegna aggiunta
    \For{$s' \in L$ in ordine inverso} \Comment{Esplora prima gli indici minori}
        \State $\operatorname{push}(S,s')$
    \EndFor
\EndWhile

\end{algorithmic}
\end{algorithm}

La funzione $\operatorname{Successori}(s)$ applica la regola di espansione
definita nella sottosezione precedente. Se non rimangono consegne da
effettuare, restituisce una lista vuota. I successori vengono inseriti
nella pila in ordine decrescente di indice, così che quello con indice
minore venga estratto per primo. La visita termina quando la pila è vuota.
L'operazione \textbf{visita} indica l'elaborazione dello stato estratto;
nel caso di una foglia, si può valutare il giro completo a essa associato.

La DFS determina l'ordine di esplorazione, ma non riduce di per sé
il numero di stati da esaminare. In assenza di potatura, il numero
complessivo di stati attraversati è $\Theta(n!)$, dove $n$ indica
il numero di consegne, escluso il deposito. Il costo temporale effettivo
dipende anche dalla rappresentazione e dalla generazione degli stati.
Questa crescita fattoriale del numero di stati
limita la scalabilità della ricerca esaustiva e motiva l'introduzione
del branch-and-bound, che consente di eliminare i rami che non possono
migliorare la migliore soluzione completa già trovata.


\subsection{Eliminazione dei rami tramite branch-and-bound}

Il branch-and-bound evita di approfondire i rami che non possono
migliorare la migliore soluzione completa già trovata. Si mantiene
il relativo percorso insieme al suo costo $G_{\mathrm{best}}$,
inizialmente posto a $+\infty$.

Quando viene raggiunto uno stato $s=(p,N,g)$ con $N=\emptyset$,
si completa il percorso tornando al deposito. Indicando con $u$
l'ultima località visitata, il costo del giro è:
\[
    G = g + \operatorname{dist}(u,d).
\]
Se $G < G_{\mathrm{best}}$, si aggiorna il costo migliore e si conserva
il giro corrispondente. Il valore $G_{\mathrm{best}}$ descrive quindi
una soluzione effettivamente disponibile, non un ottimo già noto.

Per valutare se proseguire da uno stato si utilizza un limite
inferiore $LB(s)$ al costo di qualsiasi giro che completi il suo
percorso parziale. Nella versione di base si adotta:
\[
    LB(s)=g(s),
\]
dove $g(s)$ è il costo accumulato nello stato. Questo limite è valido
perché le distanze sono non negative: aggiungere altri spostamenti
non può diminuire il costo già sostenuto.

Il controllo viene eseguito dopo l'estrazione dalla pila e prima
dell'eventuale generazione dei successori. Se:
\[
    LB(s) \geq G_{\mathrm{best}},
\]
lo stato viene scartato senza espanderlo. Qualunque completamento
avrebbe infatti un costo almeno pari a una soluzione già disponibile
e non potrebbe migliorarla. La ricerca mira a trovare un giro ottimo,
non a enumerare tutti i giri di costo ottimo.

Il limite scelto è economico da calcolare, ma poco informativo,
poiché non considera la distanza ancora necessaria per completare
il giro. La potatura può ridurre gli stati esplorati nella pratica,
ma la sua efficacia dipende dalle istanze e dalla qualità della
soluzione corrente; non si assume un miglioramento della complessità
nel caso peggiore.

\subsection{Orizzonte, frontiera attiva e condizioni di arresto}

L'\emph{orizzonte della ricerca} è il confine al quale si arresta
l'esplorazione esplicita corrente. Al termine di ciascuna iterazione,
si identifica con l'insieme $H_t$ degli stati generati ancora in attesa
nella pila $S$. Tali stati possono avere profondità diverse:
l'orizzonte non coincide quindi con un unico livello dell'albero.

La \emph{frontiera attiva} $F_t \subseteq H_t$ comprende gli stati
ancora eleggibili per ricevere ulteriore calcolo. Nella ricerca di
base tutti gli stati in attesa possono essere esaminati, pertanto:
\[
    F_t=H_t.
\]
La politica LIFO determina quale di questi stati viene estratto
successivamente. Il nodo scelto in una singola iterazione non deve
essere confuso con l'intera frontiera attiva.

La distinzione tra orizzonte e frontiera attiva consentirà,
nella versione regolata, di sospendere alcune alternative mantenendole
in memoria per una possibile riattivazione. I rami eliminati tramite
branch-and-bound, invece, vengono esclusi definitivamente, poiché
non possono migliorare la soluzione corrente.

La ricerca termina quando la pila è vuota oppure quando viene
raggiunto un budget prefissato $B$. Un'unità di budget corrisponde
all'estrazione e all'elaborazione di uno stato, anche se questo viene
scartato dal branch-and-bound o rappresenta una foglia. La generazione
dei suoi successori rientra nella stessa unità. Il numero di estrazioni
non coincide necessariamente con le espansioni effettive né garantisce
un tempo costante per passo.

Se la pila si svuota, tutte le alternative sono state esaminate o
escluse mediante un limite valido: il miglior giro trovato è quindi
ottimo. Se il budget termina con stati ancora in attesa, si restituisce
il miglior giro completo disponibile, senza garantirne l'ottimalità.
Qualora non sia stato ancora trovato alcun giro completo, tale esito
viene segnalato esplicitamente e non interpretato come assenza di
soluzioni ammissibili.
